\documentclass[a4paper, 12pt, french]{article}
\usepackage[utf8]{inputenc}
\usepackage[T1]{fontenc}
\usepackage{polyglossia}
\setmainlanguage{french}
\usepackage{lipsum}
%appeler les packages utiles

%renseigner les métadonnées titre/auteur/date

\begin{document}

% Affichez le titre généré à partir des métadonnées

% Épigraphe : utilisez la taille de texte la plus petite disponible en LaTeX,
% et alignez le bloc sur la droite de la page

% Résumé : le mot "Résumé" doit apparaître en gras avant le texte
% Le dendrobate est une petite grenouille arboricole vivant dans les forêts tropicales humides d'Amérique du sud. Vivement colorée, elle exprime ainsi sa toxicité face aux prédateurs. Ce signal visuel est une stratégie d'adaptation à son environnement très efficace : on parle d'aposématisme. Le dendrobate à tapirer possède différentes formes appelées « morphes » et la majorité d'entre eux ont une couleur à dominance jaune et noire. La sous-espèce « azureus », ou dendrobate bleu, se démarque toutefois avec une incroyable couleur bleue.

% Quatre sections non numérotées — cherchez la variante de \section qui supprime la numérotation
% Section 1 — Introduction (texte lorembarnak.com), mettre le lien dans le titre
% Le texte doit contenir un renvoi vers la figure insérée plus bas

% Figure : choisissez l'option de placement qui réserve une page entière à la figure
\begin{figure}[]

\end{figure}

% Section 2 — lorem ipsum simple
% \lipsum[1-7]

% \lipsum[1-7]

% Section 3 — même contenu, mais avec la citation Wikipedia intégrée entre deux blocs lipsum, mettre un lien dans le titre https://fr.wikipedia.org/wiki/DendrobatidaeWikipedia
% Utilisez un environnement adapté pour la citation

% \lipsum[1-4]

% [citation Wikipedia] Les Dendrobatidae sont un exemple d'organisme aposématique. Leur coloration vive annonce leur inappétence pour les prédateurs potentiels. On pense que l'aposématisme est apparu au moins quatre fois au sein de la famille des Dendrobatidae selon les arbres phylogénétiques, et ces grenouilles ont depuis subi des divergences spectaculaires - à la fois interspécifiques et intraspécifiques - dans leur coloration. Ceci est surprenant étant donné que ce type de mécanisme de défense dépend de la fréquence à laquelle ces animaux sont présents.

% \lipsum[5-7]


% Section 4 — Avertissement
% Créez un encadré avec tcolorbox : fond bleu, bordure jaune, texte blanc, centré,
% largeur 60% de la page — consultez la documentation des options de tcolorbox

% Affichez la liste des figures, puis la table des matières

\end{document}

